\section{Splitting a score gain}\label{sec:method}

\paragraph{Setup.}\label{sec:setup} Let $q_0(\cdot\mid x)$ and $q_1(\cdot\mid x)$ be the
predictive distributions of an old and a new frozen model over $K$ labels, and $S(q,y)$ a
bounded proper score, larger being better. Our default is the quadratic score
$S(q,y)=2q_y-\lVert q\rVert_2^2$, negative Brier loss up to a label-only constant; the log
score is the second case we report. Let $C=\phi(X)$ be a discrete grouping both models
observe, declared before the audit --- for relation prediction, the ordered subject--object
class pair. We call a value of $C$ a \emph{cell}. The gain vector
\begin{equation}
H_x(y)=S(q_1(\cdot\mid x),y)-S(q_0(\cdot\mid x),y),\qquad y=1,\ldots,K,
\label{eq:gain-vector}
\end{equation}
can be evaluated at a counterfactual label without rerunning either model, and that is all
the construction needs.

\paragraph{Transport.} For a cell $c$, let $\dT_c=\mathbb E[H_X(Y)\mid C=c]$ be the observed
gain and $\dP_c=\mathbb E_{H\mid c}\mathbb E_{Y'\mid c}[H_X(Y')]$ the gain when each case is
scored against an \emph{independent} label $Y'$ drawn from the same cell, and set
$\dR_c=\dT_c-\dP_c$. Transport keeps the cell and its label distribution and removes only
the pairing between a case and its own label.

\begin{theorem}[Transport and covariance identities]\label{thm:population}
For any integrable score contrast $H$,
\begin{gather}
\dT_c=\dP_c+\dR_c,\notag\\
\dR_c=\sum_{y=1}^K\operatorname{Cov}\!\left(H_X(y),\ind\{Y=y\}\mid C=c\right).
\label{eq:population-id}
\end{gather}
For the quadratic score, with $\Delta q_y(X)=q_1(y\mid X)-q_0(y\mid X)$,
\begin{equation}
\dR_c=2\sum_{y=1}^K\operatorname{Cov}\!\left(\Delta q_y(X),\ind\{Y=y\}\mid C=c\right).
\label{eq:cov-id}
\end{equation}
Consequently, if the gain vector is a function of $c$ alone, $\dR_c=0$.
\end{theorem}

\noindent Eq.~\eqref{eq:cov-id} is what the split means. The new model earns case-level
credit only insofar as its probability \emph{change} rises on the label the individual case
actually carries, relative to other cases in the same cell. Anything constant within a cell
--- shifting probability towards rare predicates for every (\textsf{man}, \textsf{surfboard})
relation alike --- lands in $\dP$, which is why we call it the group-level part. It is
\emph{not} simply ``prior fit'': under the quadratic score $\dP_c$ is the improvement in how
close the model's mean forecast sits to the cell's label frequencies \emph{minus} the
increase in the dispersion of its forecasts (\cref{prop:transported}). The covariance form is
due to Yates~\cite{yates1982external} and generalises: for every Bregman score, including the
log score, $\dR_c$ is the same covariance with $\Delta q_y$ replaced by the contrast of the
$y$-th gradient coordinate (\cref{thm:bregman}). All proofs are in the supplementary.

\paragraph{Why it matters which part a gain is in.} One cell, four cases, three labelled 1
and one labelled 2; the old model predicts $(\tfrac12,\tfrac12)$ for all four. Model A
predicts $(\tfrac34,\tfrac14)$ for every case --- the cell's frequencies. Model B predicts
$(\tfrac34,\tfrac14)$ for the three label-1 cases and $(\tfrac14,\tfrac34)$ for the label-2
case: the same distance from the old model, but pointed at each case's own label. Under the
quadratic score A gains $+\tfrac18$, all of it group-level, and B gains $+\tfrac38$ with a
case-level part of $+\tfrac12$. Now let deployment see this cell's labels in proportions
$(\tfrac14,\tfrac34)$ instead, a label shift. Reweighting each case by $p'(y)/p(y)$ sends A
to $-\tfrac38$ --- worse than the model it replaced, with not one prediction changed --- and
leaves B at $+\tfrac38$. The split said in advance which gain was exposed: a group-level gain
is contingent on the evaluation set's label frequencies matching deployment, and a case-level
gain is not. That is the operational reason to want them apart. (\Cref{app:worked} works the
example through.)

\paragraph{Estimator.} In a cell with $n_c\ge2$ cases, write $H_i(y)=H_{x_i}(y)$ and estimate
\begin{gather}
\hdT_c=\tfrac1{n_c}\textstyle\sum_iH_i(y_i),\qquad
\hdP_c=\tfrac1{n_c(n_c-1)}\textstyle\sum_{i\ne j}H_i(y_j),
\label{eq:sample-prior}\\
\hdR_c=\tbinom{n_c}{2}^{-1}\textstyle\sum_{i<j}A_{ij},
\label{eq:sample-reason}\\
A_{ij}=\tfrac12\{H_i(y_i)+H_j(y_j)-H_i(y_j)-H_j(y_i)\}.
\label{eq:kernel}
\end{gather}
These satisfy $\hdT_c=\hdP_c+\hdR_c$ exactly in every sample, and $\hdR_c$ is an unbiased
order-two U-statistic~\cite{hoeffding1948class} for $\dR_c$ (\cref{thm:sample}).
Transporting a case only to the \emph{other} cases' labels matters: fitting the cell's label
frequencies on all $n_c$ cases, the natural plug-in, shrinks every cell's case-level gain by
exactly $(n_c-1)/n_c$, halving it at $n_c=2$ (\cref{prop:attenuate}). With $m_{cy}$ the count
of label $y$ in cell $c$, case $i$'s transported score is
\begin{equation}
P_i=\frac{\sum_{y}m_{cy}H_i(y)-H_i(y_i)}{n_c-1},
\label{eq:linear}
\end{equation}
so the whole estimator costs $O(NK)$ on cached predictions, with no fitted model, held-out
fold or smoothing parameter. Cells are weighted by size; singletons cannot identify a
within-cell counterfactual and are excluded and reported.

\paragraph{Inference.}\label{sec:inference} With $\bar A_i=(n_c-1)^{-1}\sum_{j\ne i}A_{ij}$,
the H\'ajek influence contribution of case $i$ is
\begin{equation}
\widehat\psi_i=2\bar A_i-\hdR_c-\hdR.
\label{eq:influence}
\end{equation}
We sum $\widehat\psi_i$ within each image and use the one-way cluster-robust sandwich
variance, since relations from one photograph are not independent. The estimator is
asymptotically normal under bounded scores and non-dominating clusters (\cref{thm:clt}); at a
single cell, the total-gain interval \emph{is} a clustered Diebold--Mariano
test~\cite{diebold1995comparing} (\cref{cor:dm}). A randomization test that cyclically shifts
labels within every cell is a secondary, relation-level diagnostic.

\paragraph{Confidence matching.}\label{sec:matching} One property of the split must be part of
the procedure rather than a caveat. Because $\dR$ is a covariance between probability
\emph{movements} and labels, rescaling a fixed model's confidence moves score between the two
parts while changing no ranking, no top-1 decision and no information about any case
(\cref{sec:simulation} measures how much). Two released checkpoints that differ in
calibration therefore cannot be compared on their raw outputs. Wherever compared models differ
visibly in confidence we split the images at random, fit each model a single temperature by
maximum likelihood on one half, and run the split on the other half with each model's own
temperature applied. The protocol uses no label from the audited half.

\paragraph{Choosing the grouping.}\label{sec:coarsening} The split is conditional on $\phi$.
Merging cells changes it by an explicit between-cell covariance term of either sign
(\cref{prop:coarsen}), so the finest grouping that the benchmark's construction defends is
the conservative choice. A relevant variable left out of $\phi$ is the same term with the
fully adjusted cell as the fine partition, and \cref{prop:sensitivity} bounds it by a single
elicited correlation, in the spirit of the robustness value of Cinelli and
Hazlett~\cite{cinelli2020making}.
